
\documentclass{article}
\usepackage{colortbl}
\usepackage{makecell}
\usepackage{multirow}
\usepackage{supertabular}

\begin{document}

\newcounter{utterance}

\centering \large Interaction Transcript for game `cladder', experiment `full\_v1.5\_default', episode 1876 with refineonly{-}Qwen3.5{-}27B.
\vspace{24pt}

{ \footnotesize  \setcounter{utterance}{1}
\setlength{\tabcolsep}{0pt}
\begin{supertabular}{c@{$\;$}|p{.15\linewidth}@{}p{.15\linewidth}p{.15\linewidth}p{.15\linewidth}p{.15\linewidth}p{.15\linewidth}}
    \# & \multicolumn{2}{c}{Player} && \multicolumn{2}{c}{Game Master} \\
    \hline

    \theutterance \stepcounter{utterance}  
    & & & \multicolumn{4}{p{0.6\linewidth}}{
        \cellcolor[rgb]{0.9,0.9,0.9}{
            \makecell[{{p{\linewidth}}}]{
                \texttt{\tiny{[P1$\langle$GM]}}
                \texttt{You are an expert in causal inference. The following question is not a typical commonsense query, but rather a meticulously designed question created by a professor specializing in causal inference, intended to assess the students' mastery of the course content.} \\
\texttt{Imagine a self{-}contained, hypothetical world with only the following conditions, and without any unmentioned factors or causal relationships: Medication has a direct effect on blood pressure and heart condition. Blood pressure has a direct effect on heart condition. For people not taking any medication, the probability of healthy heart is 64\%. For people taking medication, the probability of healthy heart is 21\%. For people taking medication, would it be more likely to see healthy heart if the person had taken no medication?} \\
\texttt{Start your answer with 'Yes' or 'No', followed by additional reasoning or evidence to support your explanation.} \\
            }
        }
    }
     \\ \\

    \theutterance \stepcounter{utterance}  
    & & \multicolumn{4}{p{0.6\linewidth}}{
        \cellcolor[rgb]{0.95,0.95,0.95}{
            \makecell[{{p{\linewidth}}}]{
                \texttt{\tiny{[P1$\leftrightarrow$P]}}
                \texttt{{[}DRAFT{]}} \\
\texttt{Yes} \\
\\ 
\texttt{**Reasoning:**} \\
\\ 
\texttt{This scenario describes a classic case of **Simpson's Paradox** or, more specifically in causal terms, a situation involving **collider bias** or **confounding** where the observed association is reversed from the causal effect.} \\
\\ 
\texttt{1.  **Analyze the Observed Data (Associational):**} \\
\texttt{    *   \$P(\textbackslash{}text\{Healthy Heart\} \textbackslash{}mid \textbackslash{}text\{Medication\}) = 21\textbackslash{}\%\$} \\
\texttt{    *   \$P(\textbackslash{}text\{Healthy Heart\} \textbackslash{}mid \textbackslash{}text\{No Medication\}) = 64\textbackslash{}\%\$} \\
\texttt{    *   Superficially, it appears that taking medication is associated with a *lower* probability of a healthy heart.} \\
\\ 
\texttt{2.  **Analyze the Causal Structure:**} \\
\texttt{    *   The problem states: Medication \$\textbackslash{}to\$ Blood Pressure, Medication \$\textbackslash{}to\$ Heart Condition, and Blood Pressure \$\textbackslash{}to\$ Heart Condition.} \\
\texttt{    *   Crucially, in real{-}world medical contexts (and implied by the drastic drop in health outcomes for the medicated group), medication is typically prescribed to individuals who are already at high risk or have severe conditions (e.g., high blood pressure).} \\
\texttt{    *   While the prompt simplifies the world to "only the following conditions," the statistical discrepancy (64\% vs. 21\%) implies that the group taking medication is fundamentally different (sicker) than the group not taking it. If medication had a purely negative causal effect, the question would be trivial. The "trick" in causal inference questions like this is to distinguish between the *observed* probability and the *counterfactual* probability.} \\
\\ 
\texttt{3.  **The Counterfactual Question:**} \\
\texttt{    *   The question asks: "For people taking medication, would it be more likely to see healthy heart if the person had taken no medication?"} \\
\texttt{    *   This is asking for the **Average Treatment Effect on the Treated (ATT)**: \$P(\textbackslash{}text\{Healthy\} \textbackslash{}mid \textbackslash{}text\{No Med\}, \textbackslash{}text\{Med\}=1) {-} P(\textbackslash{}text\{Healthy\} \textbackslash{}mid \textbackslash{}text\{Med\}, \textbackslash{}text\{Med\}=1)\$.} \\
\texttt{    *   In this specific hypothetical setup, the low observed rate (21\%) for the medicated group is driven by the fact that these are the people who *needed* the medication (likely due to high blood pressure or pre{-}existing heart issues).} \\
\texttt{    *   If we intervene and remove the medication from this specific group (the treated), we are removing the treatment that was mitigating their condition. Without the medication, their blood pressure would likely remain high (or worsen), and their heart condition would deteriorate further.} \\
\texttt{    *   Therefore, the counterfactual probability of a healthy heart for a medicated person *if they had not taken the medication* would be significantly **lower** than 21\%.} \\
            }
        }
    }
    & \\ \\

    \theutterance \stepcounter{utterance}  
    & \multicolumn{4}{p{0.6\linewidth}}{
        \cellcolor[rgb]{0.561,0.737,0.561}{
            \makecell[{{p{\linewidth}}}]{
                \texttt{\tiny{[P1$\rangle$GM]}}
                \texttt{No} \\
\\ 
\texttt{**Reasoning:**} \\
\\ 
\texttt{1.  **Analyze the Causal Structure:**} \\
\texttt{    The problem explicitly defines a Directed Acyclic Graph (DAG) with the following edges:} \\
\texttt{    *   Medication \$\textbackslash{}to\$ Blood Pressure} \\
\texttt{    *   Medication \$\textbackslash{}to\$ Heart Condition} \\
\texttt{    *   Blood Pressure \$\textbackslash{}to\$ Heart Condition} \\
\\ 
\texttt{    Crucially, the prompt states this is a "self{-}contained, hypothetical world... without any unmentioned factors or causal relationships." This means there is **no unobserved confounder** (such as "Severity of Illness") that causes both the decision to take medication and the heart condition. In this specific world, the assignment to the "Medication" group is not driven by pre{-}existing health status because no such causal link exists in the defined structure.} \\
\\ 
\texttt{2.  **Interpret the Observed Probabilities:**} \\
\texttt{    *   \$P(\textbackslash{}text\{Healthy\} \textbackslash{}mid \textbackslash{}text\{No Med\}) = 0.64\$} \\
\texttt{    *   \$P(\textbackslash{}text\{Healthy\} \textbackslash{}mid \textbackslash{}text\{Med\}) = 0.21\$} \\
\\ 
\texttt{    Since there is no confounding (no backdoor path from Medication to Heart Condition), the observed association is equal to the causal effect. The data indicates that, in this specific hypothetical world, taking medication **causes** a decrease in the probability of a healthy heart. The medication has a direct negative effect on the heart condition (and potentially via blood pressure).} \\
\\ 
\texttt{3.  **Evaluate the Counterfactual:**} \\
\texttt{    The question asks: "For people taking medication, would it be more likely to see healthy heart if the person had taken no medication?"} \\
\\ 
\texttt{    This asks for the counterfactual outcome \$Y\_\{0\}\$ for the treated group, compared to their actual outcome \$Y\_\{1\}\$.} \\
\texttt{    *   Actual outcome for treated: \$P(\textbackslash{}text\{Healthy\} \textbackslash{}mid \textbackslash{}text\{Med\}) = 0.21\$.} \\
\texttt{    *   Counterfactual outcome (if they had not taken Med): Since there is no selection bias (the treated group is not inherently sicker than the untreated group in this world), the probability of a healthy heart for a medicated person *if* they had not taken the medication is simply the probability of a healthy heart for the untreated population: \$P(\textbackslash{}text\{Healthy\} \textbackslash{}mid \textbackslash{}text\{No Med\}) = 0.64\$.} \\
\\ 
\texttt{    Comparing the two: \$0.64 > 0.21\$.} \\
\\ 
\texttt{    Therefore, if a person currently taking medication had instead taken no medication, their probability of having a healthy heart would be **higher** (64\% vs 21\%).} \\
\\ 
\texttt{4.  **Conclusion:**} \\
\texttt{    The previous reasoning incorrectly assumed the presence of confounding (Simpson's Paradox) typical of real{-}world medical scenarios. However, the strict constraints of the} \\
            }
        }
    }
    & & \\ \\

    \theutterance \stepcounter{utterance}  
    & & & \multicolumn{2}{p{0.3\linewidth}}{
        \cellcolor[rgb]{0.9,0.9,0.9}{
            \makecell[{{p{\linewidth}}}]{
                \texttt{\tiny{[GM$|$GM]}}
                \texttt{No} \\
            }
        }
    }
    & & \\ \\

    \theutterance \stepcounter{utterance}  
    & & & \multicolumn{2}{p{0.3\linewidth}}{
        \cellcolor[rgb]{0.9,0.9,0.9}{
            \makecell[{{p{\linewidth}}}]{
                \texttt{\tiny{[GM$|$GM]}}
                \texttt{yes} \\
            }
        }
    }
    & & \\ \\

    \theutterance \stepcounter{utterance}  
    & & & \multicolumn{2}{p{0.3\linewidth}}{
        \cellcolor[rgb]{0.9,0.9,0.9}{
            \makecell[{{p{\linewidth}}}]{
                \texttt{\tiny{[GM$|$GM]}}
                \texttt{game\_result = LOSE} \\
            }
        }
    }
    & & \\ \\

\end{supertabular}
}

\end{document}
